\documentclass[12pt]{article}
\usepackage{amsmath, amssymb, amsthm}
\usepackage[margin=1in]{geometry}

\newtheorem{theorem}{Theorem}
\newtheorem{definition}[theorem]{Definition}
\newtheorem{lemma}[theorem]{Lemma}

\newcommand{\next}[1]{#1_{\text{next}}}
\newcommand{\falling}[2]{{#1}^{\underline{#2}}}

\title{Discrete Calculus: From the Line to the Circle}
\author{An Algebra II Journey}
\date{}

\begin{document}
\maketitle

\section*{Introduction}
In ordinary algebra, you learn to analyze how quantities change by looking at their rates. Calculus is the mathematics of change, but traditional calculus relies on limits and continuous motion. Discrete calculus achieves the same structural beauty by looking at steps rather than slides. 

Because we are stepping, we will strictly avoid the continuous shift operator $x+1$. Instead, we will build our entire system using only the current state of a function, $f$, and its discrete difference, $\Delta f$. Our guiding algebraic engine will be the simple but profound identity:
\[ \next{f} = f + \Delta f \]
This tells us that the ``next'' value of any function is just its current value plus its difference. 

\part{The Flat Line}

\section{The Forward Difference}

\begin{definition}
The \textbf{forward difference} of a function $f$, written $\Delta f$, is the change in the function from its current value to its next value:
\[ \Delta f = \next{f} - f \]
\end{definition}

Because $\next{f} = f + \Delta f$, we can understand how functions step forward purely algebraically.

\begin{theorem}[Linearity]
For any functions $f, g$ and constant $c$:
\begin{enumerate}
    \item $\Delta(cf) = c\Delta f$
    \item $\Delta(f + g) = \Delta f + \Delta g$
\end{enumerate}
\end{theorem}

\begin{proof}
\begin{enumerate}
    \item $\next{(cf)} - cf = c(\next{f}) - cf = c(\next{f} - f) = c\Delta f$.
    \item $\next{(f+g)} - (f+g) = (\next{f} + \next{g}) - (f+g) = (\next{f} - f) + (\next{g} - g) = \Delta f + \Delta g$. \qedhere
\end{enumerate}
\end{proof}

Let us compute the differences of basic powers. Since $\Delta x = 1$, we have $\next{x} = x + 1$.
\begin{align*}
    \Delta x &= 1 \\
    \Delta x^2 &= \next{(x^2)} - x^2 = (\next{x})^2 - x^2 = (x+1)^2 - x^2 = 2x + 1 \\
    \Delta x^3 &= \next{(x^3)} - x^3 = (\next{x})^3 - x^3 = (x+1)^3 - x^3 = 3x^2 + 3x + 1
\end{align*}
These formulas are messy. They do not mirror the elegant power rule of continuous calculus. To find an elegant power rule, we must change our representation.

\section{Falling Powers (The Permutation Function)}

In Algebra II, you learned the permutation formula $P(n, k) = n(n-1)(n-2)\dots(n-k+1)$, which counts the number of ways to arrange $k$ items from $n$ choices. We define our discrete powers using this exact structure.

\begin{definition}
The \textbf{falling power} $\falling{x}{k}$ (read ``$x$ to the $k$ falling'') is defined as:
\[ \falling{x}{k} = x(x-1)(x-2)\dots(x-k+1) \]
For $k=0$, we define $\falling{x}{0} = 1$.
\end{definition}

Notice that this is just $P(x, k)$. For example:
\[ \falling{x}{2} = x(x-1) = x^2 - x \]
\[ \falling{x}{3} = x(x-1)(x-2) = x^3 - 3x^2 + 2x \]
Falling powers are simply expanded polynomials. They are a different \textit{basis} for writing polynomials, just as you can write $10$ as $10^1$ or $2^3 + 2^1$.

\section{The Discrete Product Rule}

Before proving the power rule for falling powers, we need a product rule. In continuous calculus, the product rule is $\frac{d}{dx}(fg) = f\frac{dg}{dx} + g\frac{df}{dx}$. In discrete calculus, stepping introduces an interaction term.

\begin{theorem}[Discrete Product Rule]
\[ \Delta(fg) = f\Delta g + g\Delta f + \Delta f \Delta g \]
\end{theorem}

\begin{proof}
We use our engine $\next{f} = f + \Delta f$.
\begin{align*}
    \Delta(fg) &= \next{(fg)} - fg \\
    &= \next{f}\next{g} - fg \\
    &= (f + \Delta f)(g + \Delta g) - fg \\
    &= fg + f\Delta g + g\Delta f + \Delta f \Delta g - fg \\
    &= f\Delta g + g\Delta f + \Delta f \Delta g
\end{align*}
Notice how the $\Delta f \Delta g$ term arises naturally from the FOIL expansion, without ever needing to write $f(x+1)$. 
\end{proof}

\section{The Discrete Power Rule}

Now we can prove that falling powers have the elegant difference rule we wanted.

\begin{theorem}[Discrete Power Rule]
\[ \Delta(\falling{x}{k}) = k \falling{x}{k-1} \]
\end{theorem}

\begin{proof}
We proceed by induction on $k$.

\textbf{Base Case ($k=1$):} 
$\falling{x}{1} = x$. We know $\Delta x = 1$. The formula gives $1 \falling{x}{0} = 1(1) = 1$. The base case holds.

\textbf{Inductive Step:} 
Assume $\Delta(\falling{x}{k-1}) = (k-1)\falling{x}{k-2}$, and write $F = \falling{x}{k-2}$. Then $\falling{x}{k-1} = (x-k+2)F$.
We can write $\falling{x}{k} = \falling{x}{k-1} \cdot (x - k + 1)$.
Let $f = \falling{x}{k-1}$ and $g = x - k + 1$.
Notice $\Delta g = \Delta(x - k + 1) = \Delta x - \Delta k + \Delta 1 = 1 - 0 + 0 = 1$.
Applying the Product Rule:
\begin{align*}
    \Delta(\falling{x}{k}) &= g \Delta f + f \Delta g + \Delta f \Delta g \\
    &= (x - k + 1)(k-1)F + (x-k+2)F + (k-1)F
\end{align*}
We factor out $F$:
\[ F \big[ (k-1)(x-k+1) + (x-k+2) + (k-1) \big] = F \big[ (k-1)(x-k+1) + x + 1 \big] \]
Simplify the bracket by expanding $(k-1)(x-k+1)$:
\begin{align*}
    (k-1)x - (k-1)^2 + x + 1 &= kx - (k^2 - 2k + 1) + 1 \\
    &= kx - k^2 + 2k \\
    &= k(x - k + 2)
\end{align*}
Substitute this back into our factored expression:
\[ F \big[ k(x - k + 2) \big] = k(x-k+2)F = k \falling{x}{k-1} \]
This completes the induction.
\end{proof}

\section{The Discrete Reciprocal and Quotient Rules}

\begin{theorem}[Reciprocal Rule]
\[ \Delta\left(\frac{1}{f}\right) = \frac{-\Delta f}{f(f + \Delta f)} \]
\end{theorem}

\begin{proof}
Let $f \cdot \frac{1}{f} = 1$. The difference of a constant is 0, so:
\[ \Delta\left(f \cdot \frac{1}{f}\right) = 0 \]
Apply the Product Rule, with $g = 1/f$:
\[ f \Delta\left(\frac{1}{f}\right) + \frac{1}{f}\Delta f + \Delta f \Delta\left(\frac{1}{f}\right) = 0 \]
Group the $\Delta(1/f)$ terms:
\[ \Delta\left(\frac{1}{f}\right) (f + \Delta f) = -\frac{\Delta f}{f} \]
Divide by $(f + \Delta f)$, which is exactly $\next{f}$:
\[ \Delta\left(\frac{1}{f}\right) = \frac{-\Delta f}{f(f + \Delta f)} \]
\end{proof}

\begin{theorem}[Quotient Rule]
\[ \Delta\left(\frac{f}{g}\right) = \frac{g\Delta f - f\Delta g}{g(g + \Delta g)} \]
\end{theorem}

\begin{proof}
Let $h = f \cdot \frac{1}{g}$. Apply the Product Rule:
\[ \Delta h = f\Delta\left(\frac{1}{g}\right) + \frac{1}{g}\Delta f + \Delta f \Delta\left(\frac{1}{g}\right) \]
Substitute the Reciprocal Rule for $\Delta(1/g)$:
\[ \Delta h = f\left(\frac{-\Delta g}{g(g + \Delta g)}\right) + \frac{\Delta f}{g} + \Delta f\left(\frac{-\Delta g}{g(g + \Delta g)}\right) \]
Factor out $\frac{1}{g(g + \Delta g)}$:
\[ \Delta h = \frac{1}{g(g + \Delta g)} \Big[ -f\Delta g + \Delta f(g + \Delta g) - \Delta f \Delta g \Big] \]
Expand and cancel the $\Delta f \Delta g$ terms:
\[ \Delta h = \frac{1}{g(g + \Delta g)} \Big[ -f\Delta g + g\Delta f + \Delta f \Delta g - \Delta f \Delta g \Big] = \frac{g\Delta f - f\Delta g}{g(g + \Delta g)} \]
\end{proof}

\section{The Exponential Base 2}

\begin{theorem}
For the function $2^n$:
\[ \Delta(2^n) = 2^n \]
\end{theorem}

\begin{proof}
By the laws of exponents, doubling a power of 2 shifts it to the next power:
\[ \next{(2^n)} = 2 \cdot 2^n \]
But by our engine, $\next{(2^n)} = 2^n + \Delta(2^n)$. Equate the two:
\[ 2^n + \Delta(2^n) = 2 \cdot 2^n \]
Subtract $2^n$ from both sides:
\[ \Delta(2^n) = 2 \cdot 2^n - 2^n = 2^n \]
\end{proof}

\section{Summation and the Fundamental Theorem}

If $\Delta$ is our discrete derivative, then summation $\sum$ is our discrete integral.

\begin{theorem}[Fundamental Theorem of Discrete Calculus]
For any function $f$, summing its difference yields the original function evaluated at the boundaries:
\[ \sum_{x=m}^{n-1} \Delta f(x) = f(n) - f(m) \]
\end{theorem}

\begin{proof}
Recall that $\Delta f(x) = \next{f}(x) - f(x)$. If we imagine $x$ taking integer steps from $m$ to $n-1$, the ``next'' value of $f$ at any step $x$ is simply $f$ evaluated at $x+1$. We write this as $f(x+1)$. 
\begin{align*}
    \sum_{x=m}^{n-1} \Delta f(x) &= \sum_{x=m}^{n-1} \big( f(x+1) - f(x) \big) \\
    &= \big(f(m+1) - f(m)\big) + \big(f(m+2) - f(m+1)\big) + \dots + \big(f(n) - f(n-1)\big)
\end{align*}
Every interior term cancels out (telescoping), leaving only the final boundary and the initial boundary:
\[ f(n) - f(m) \]
\end{proof}

This gives us immense power. For example, to find $\sum_{x=0}^{n-1} x^2$, we traditionally need binomial expansion. But we know $x^2 = \falling{x}{2} + \falling{x}{1}$, and $\Delta(\falling{x}{3}) = 3\falling{x}{2}$, so $\falling{x}{2} = \Delta\left(\frac{1}{3}\falling{x}{3}\right)$. 
\[ \sum_{x=0}^{n-1} x^2 = \sum_{x=0}^{n-1} \left( \Delta\left(\frac{1}{3}\falling{x}{3}\right) + \Delta\left(\frac{1}{2}\falling{x}{2}\right) \right) = \left[ \frac{1}{3}\falling{x}{3} + \frac{1}{2}\falling{x}{2} \right]_0^n = \frac{n(n-1)(n-2)}{3} + \frac{n(n-1)}{2} \]

\section{Summation by Parts}

\begin{theorem}[Summation by Parts]
\[ \sum_{x=m}^{n-1} f(x) \Delta g(x) = \Big[ f(x)g(x) \Big]_m^n - \sum_{x=m}^{n-1} g(x) \Delta f(x) - \sum_{x=m}^{n-1} \Delta f(x) \Delta g(x) \]
\end{theorem}

\begin{proof}
Start with the Product Rule:
\[ f \Delta g = \Delta(fg) - g \Delta f - \Delta f \Delta g \]
Sum both sides from $m$ to $n-1$:
\[ \sum_{x=m}^{n-1} f \Delta g = \sum_{x=m}^{n-1} \Delta(fg) - \sum_{x=m}^{n-1} g \Delta f - \sum_{x=m}^{n-1} \Delta f \Delta g \]
Apply the Fundamental Theorem to the first term on the right:
\[ \sum_{x=m}^{n-1} \Delta(fg) = f(n)g(n) - f(m)g(m) = \Big[ fg \Big]_m^n \]
Substituting this back completes the proof.
\end{proof}

\section{The Harmonic Series}

In continuous calculus, the integral of $1/x$ is $\ln(x)$. Discretely, we define the Harmonic series:
\[ H_n = \sum_{k=1}^n \frac{1}{k} = 1 + \frac{1}{2} + \frac{1}{3} + \dots + \frac{1}{n} \]
Its difference is exactly $\Delta H_n = H_{n+1} - H_n = \frac{1}{n+1}$. Let us prove a crucial property of this series: it grows without bound.

\begin{theorem}
The Harmonic series diverges. That is, for any number $M$, there exists an $n$ such that $H_n > M$.
\end{theorem}

\begin{proof}
We group the terms of the series. Each group is twice as long as the previous one:
\begin{align*}
    H_1 &= 1 \\
    H_2 &= 1 + \frac{1}{2} \\
    H_4 &= 1 + \frac{1}{2} + \left(\frac{1}{3} + \frac{1}{4}\right) \\
    H_8 &= 1 + \frac{1}{2} + \left(\frac{1}{3} + \frac{1}{4}\right) + \left(\frac{1}{5} + \frac{1}{6} + \frac{1}{7} + \frac{1}{8}\right)
\end{align*}
Notice that $\frac{1}{3} + \frac{1}{4} > \frac{1}{4} + \frac{1}{4} = \frac{1}{2}$.
Similarly, $\frac{1}{5} + \frac{1}{6} + \frac{1}{7} + \frac{1}{8} > \frac{1}{8} + \frac{1}{8} + \frac{1}{8} + \frac{1}{8} = \frac{1}{2}$.
Every grouped block after the first two terms is strictly greater than $\frac{1}{2}$, and the first two terms already sum to $1+\frac12$. 
Therefore, $H_{2^k} \ge 1 + k\left(\frac{1}{2}\right)$. Since $k$ can be made as large as we want, $H_n$ eventually exceeds any number $M$.
\end{proof}

\section{Stirling Numbers of the Second Kind}

We have established that falling powers have a simple difference. But how do we find the difference of standard powers like $\Delta x^3$? We change the basis using Stirling Numbers.

\begin{theorem}[Stirling Recurrence]
\[ x \cdot \falling{x}{k} = \falling{x}{k+1} + k \falling{x}{k} \]
\end{theorem}

\begin{proof}
Factor $\falling{x}{k}$ out of the right side:
\begin{align*}
    \falling{x}{k+1} + k \falling{x}{k} &= (x-k)\falling{x}{k} + k \falling{x}{k} \\
    &= \falling{x}{k}(x - k + k) \\
    &= x \falling{x}{k}
\end{align*}
\end{proof}

This recurrence allows us to express standard powers in terms of falling powers. For example:
\begin{align*}
    x^1 &= \falling{x}{1} \\
    x^2 &= x \falling{x}{1} = \falling{x}{2} + 1\falling{x}{1} = \falling{x}{2} + \falling{x}{1} \\
    x^3 &= x \falling{x}{2} + x \falling{x}{1} = (\falling{x}{3} + 2\falling{x}{2}) + (\falling{x}{2} + \falling{x}{1}) = \falling{x}{3} + 3\falling{x}{2} + \falling{x}{1}
\end{align*}
The coefficients of these expansions $(1), (1, 1), (1, 3, 1)$ are the Stirling Numbers of the Second Kind. By rewriting standard powers in the falling basis, their differences become trivial to compute using the Power Rule.

\part{The Curved Circle}

\section{Rational Parametrization of the Circle}

We now develop a genuinely discrete analogue of trigonometry, entirely algebraic, without angles or limits. Start with the rational parametrization of the unit circle:
\[ C(t) = \frac{1-t^2}{1+t^2}, \qquad S(t) = \frac{2t}{1+t^2} \]

\begin{theorem}
$C(t)$ and $S(t)$ satisfy the Pythagorean identity:
\[ C(t)^2 + S(t)^2 = 1 \]
\end{theorem}

\begin{proof}
\[ C(t)^2 + S(t)^2 = \left(\frac{1-t^2}{1+t^2}\right)^2 + \left(\frac{2t}{1+t^2}\right)^2 = \frac{1 - 2t^2 + t^4 + 4t^2}{(1+t^2)^2} = \frac{1 + 2t^2 + t^4}{(1+t^2)^2} = \frac{(1+t^2)^2}{(1+t^2)^2} = 1 \]
\end{proof}

\section{The $\oplus$ Operation}

To add angles discretely, we define the $\oplus$ operation on the parameter $t$:
\[ t \oplus u = \frac{t+u}{1-tu} \]

\begin{lemma}
$\oplus$ behaves like addition:
\begin{enumerate}
    \item \textbf{Identity:} $t \oplus 0 = t$
    \item \textbf{Inverse:} $t \oplus (-t) = 0$
    \item \textbf{Commutativity:} $t \oplus u = u \oplus t$
\end{enumerate}
\end{lemma}

\begin{proof}
\begin{enumerate}
    \item $\frac{t+0}{1-t(0)} = \frac{t}{1} = t$
    \item $\frac{t+(-t)}{1-t(-t)} = \frac{0}{1+t^2} = 0$
    \item $\frac{t+u}{1-tu} = \frac{u+t}{1-ut}$ \qedhere
\end{enumerate}
\end{proof}

\section{Discrete Circular Addition Formulas}

The most vital property of our discrete functions is that they obey addition formulas strictly through rational algebra.

\begin{theorem}
\[ C(t \oplus u) = C(t)C(u) - S(t)S(u) \]
\[ S(t \oplus u) = S(t)C(u) + C(t)S(u) \]
\end{theorem}

\begin{proof}
We prove the $C$ identity; the $S$ identity follows similarly. Substitute $t \oplus u$ into $C$:
\[ C(t \oplus u) = \frac{1 - \left(\frac{t+u}{1-tu}\right)^2}{1 + \left(\frac{t+u}{1-tu}\right)^2} \]
Multiply numerator and denominator by $(1-tu)^2$:
\[ = \frac{(1-tu)^2 - (t+u)^2}{(1-tu)^2 + (t+u)^2} \]
Expand both squares:
\[ = \frac{(1 - 2tu + t^2u^2) - (t^2 + 2tu + u^2)}{(1 - 2tu + t^2u^2) + (t^2 + 2tu + u^2)} = \frac{1 - 4tu + t^2u^2 - t^2 - u^2}{1 + t^2u^2 + t^2 + u^2} \]
Factor the numerator and denominator by grouping:
\[ = \frac{(1-t^2)(1-u^2) - (2t)(2u)}{(1+t^2)(1+u^2)} \]
Split the fraction:
\[ = \frac{1-t^2}{1+t^2}\frac{1-u^2}{1+u^2} - \frac{2t}{1+t^2}\frac{2u}{1+u^2} = C(t)C(u) - S(t)S(u) \]
\end{proof}

\section{The Circular Difference Operator}

We define a difference operator tailored to stepping along the circle:

\begin{definition}
The \textbf{circular difference} of $f$ with step $u$ is:
\[ \Delta_{\oplus, u} f(t) = f(t \oplus u) - f(t) \]
\end{definition}

Just as our flat engine was $\next{f} = f + \Delta f$, our circular engine is $f(t \oplus u) = f(t) + \Delta_{\oplus, u} f(t)$.

\begin{theorem}[Discrete Circular Derivatives]
\begin{align*}
    \Delta_{\oplus, u} C(t) &= (C(u)-1)C(t) - S(u)S(t) \\
    \Delta_{\oplus, u} S(t) &= S(u)C(t) + (C(u)-1)S(t)
\end{align*}
\end{theorem}

\begin{proof}
Using the addition formulas:
\begin{align*}
    \Delta_{\oplus, u} C(t) &= C(t \oplus u) - C(t) \\
    &= \big( C(t)C(u) - S(t)S(u) \big) - C(t) \\
    &= C(t)C(u) - C(t) - S(t)S(u) \\
    &= (C(u)-1)C(t) - S(u)S(t)
\end{align*}
The proof for $S(t)$ is identical.
\end{proof}

\section{Summation on the Circle}

The rational parametrization immediately gives us:
\[ 1 + C(t) = 1 + \frac{1-t^2}{1+t^2} = \frac{2}{1+t^2} \implies \frac{1}{1+t^2} = \frac{1+C(t)}{2} \]
This is the discrete analogue of the integral $\int \frac{1}{1+t^2} dt = \arctan(t)$. 

If we have a sequence of parameters $t_0, t_1, t_2 \dots t_N$ stepping by $u$ (so $t_{n} = t_{n-1} \oplus u$), we can sum the circular difference using telescoping, identical to the flat FTC:
\[ \sum_{n=0}^{N-1} \Delta_{\oplus, u} f(t_n) = \sum_{n=0}^{N-1} \big( f(t_n \oplus u) - f(t_n) \big) = f(t_N) - f(t_0) \]

\part{The Grand Principle}

\section{The Failure of the Chain Rule}

In ordinary calculus, the chain rule makes substitution powerful: $\frac{d}{dx}(f(g(x))) = f'(g(x))g'(x)$. 
In discrete calculus, there is no direct chain rule. Applying a function to a sequence $g(n)$ yields:
\[ \Delta(f \circ g)(n) = f(\next{g}(n)) - f(g(n)) \]
This does not factor into $(\Delta f)(g(n)) \cdot \Delta g(n)$. The product rule itself exhibits this irreducible interaction via the $\Delta f \Delta g$ term.

\section{The Art of Changing Representation}

Because substitution fails, the natural discrete strategy is not \textit{change of variables}, but \textbf{change of representation}. When a sum does not telescope, we ask: what representation makes its finite difference visible?

\begin{itemize}
    \item \textbf{Flat powers fail:} $\Delta x^k$ is a messy polynomial.
    \item \textbf{Falling powers win:} $\Delta \falling{x}{k} = k \falling{x}{k-1}$. We changed the basis from $x^k$ to $\falling{x}{k}$ to fit the operator $\Delta$.
    \item \textbf{Flat rationals fail:} $\Delta \left(\frac{1}{1+t^2}\right)$ is intractable under standard steps.
    \item \textbf{Circular functions win:} $\frac{1}{1+t^2} = \frac{1+C(t)}{2}$. We changed the representation to fit the operator $\Delta_{\oplus}$.
\end{itemize}

\begin{quote}
\textit{Discrete calculus is substantially the art of finding the basis or representation adapted to the operator.}
\end{quote}

\end{document}